\scanpage{126}
\mspage{\handwritten{2}}

de schémas en groupes finis et plats sur $S$
\[
e\handwritten{\to}H_{\mathrm{inf}}\handwritten{\to}H\handwritten{\to}H_{\text{sép}}\handwritten{\to}e
\]
où $H_{\mathrm{inf}}$ est purement infinitésimal (i.e. la projection $H_{\mathrm{inf}}\handwritten{\to}S$
est géométriquement injectif) et où $H_{\text{sép}}$ est non ramifié sur $S$.
(Dans le cas d'un corps de base $k$, une telle suite exacte est
déduite d'une suite exacte analogue pour un groupe algébrique
fini \inserted{$G$} sur $k$; d'ailleurs si $k$ est parfait, cette suite splitte canoniquement, car $G_{\text{réd}}$ est \inserted{alors} un sous-groupe de $G$, isomorphe à $G_{\text{sép}}$ par la
projection $G\quad G_{\text{sép}}$. Faire attention cependant que $G_{\text{réd}}$ va \unsure{opérer|opérser}
sur $G_{\mathrm{inf}}$, on aura seulement un produit semi-direct). Enfin

\deleted{\unsure{La condition (v) est certainement provisoire}}

(vi) $S$ est réduit.

Les conditions (v) et (vi) ont une sale gueule, et sont essentiellement provisoires. Elles servent à \deleted{assurer la validité} \inserted{\handwritten{cause}} du

\uline{Lemme} Soit $H$ comme dans (v), $S$ comme dans (vi), $P$ un fibré principal homogène sous $H$, $Q$ un autre, $u,v$ deux isomorphismes de $P$ dans $Q$
transformant un \handwritten{``}point \deleted{\unsure{rationnel}} \inserted{marqué}\handwritten{''} donné en un autre donné, alors
$u=v$. (On est ramené, en tordant $H$, au cas où $P$ est trivial, et
alors cela résulte du fait suivant: une section de $Q$ est un isomorphisme de $S$ sur une composante connexe de $Q_{\text{réd}}$).

Remarque néanmoins que les conditions (i) à (vi) n'excluent
pas des groupes structuraux tordus (relativement à un corps de base
$k$).

Soit maintenant $a$ un ``point marqué'' de $S$ (i.e. une extension
algébriquement close du corps résiduel d'un $s\handwritten{\in}S$). Pour tout $G\handwritten{\in}\underline G$,
on désigne par $Z(S,a;G)$ ou simplement $Z(G)$ l'ensemble des classes
(à isomorphisme près) de fibrés principaux homogènes sur $S$, de
groupe $F(G)$, munis d'un point marqué au dessus de $a$. Evidemment
$Z(G)$ est un foncteur de $\underline G$ dans la catégorie des ensembles. Ce
foncteur satisfait aux conditions suivantes: 1) Il commute aux
produits (\deleted{\unsure{trivial), et|trivial}} car $F$ y commute) 2) il commute aux noyaux
de couples, en d'autres termes si $G''\handwritten{\to}G\handwritten{\rightrightarrows}G'$ est exact dans $\underline G$,
alors $Z(G'')\handwritten{\to}Z(G)\handwritten{\rightrightarrows}Z(G')$ est exacte, i.e. $Z(G'')$ s'identifie
à l'ensemble des éléments de $Z(G)$ dont l'image dans $Z(G')$ par $Z(u)$
\deleted{\unsure{(}}et $Z(v\deleted{\unsure{a}})$ est la même. (Ceci résulte de l'exactitu\unsure{de|é} de $F$, et du
lemme). \unsure{t|c}
% archive sheet number: 125
